\chapter{The Emergence Cocycle}\label{ch:emergence}

\epigraph{\itshape ``The whole is something besides the parts.''}{Aristotle, \emph{Metaphysics} $\Delta$, 1041b.}

The previous chapters built the kinematics of an idea algebra: a monoid
$\IdMon$, a resonance pairing $\rho$, and the Aufhebung decomposition
$a\compose b = \sigma(a,b) + \nu(a,b) + \eta(a,b)$. These data describe how
ideas combine. They do not yet describe what is \emph{lost} or \emph{gained}
when they do. The slogan of the present chapter is that whatever quantitative
content is carried by an idea --- its cost, its information, its predictive
accuracy, its survival fitness, its degree of belief --- almost never
distributes over composition. The combined idea $a\compose b$ is, as a
quantity, almost never the sum of $a$ and $b$. The discrepancy is a function
of the pair $(a,b)$, and the central observation of the chapter is that this
function, no matter what valuation it is the discrepancy of, satisfies a
single algebraic constraint forced by associativity. That constraint is the
2-cocycle identity, and the discrepancy is the \emph{emergence cocycle}
$\omega$.

We will derive $\omega$ from first principles, prove the cocycle identity
(Theorem~\ref{thm:cocycle}), normalize $\omega$ at the unit, study its
indeterminacy under coboundaries (Proposition~\ref{prop:cob}), and prove the
Transmission Telescoping Theorem (Theorem~\ref{thm:telescope}) which
expresses the value of an iterated composition as the sum of its parts plus a
sum of pairwise emergence terms. We then make precise the distinction
between \emph{weak} and \emph{strong} emergence in cohomological language:
weak emergence corresponds to a uniformly bounded $\omega$, strong emergence
to an $\omega$ that is non-trivial in the second cohomology
$H^2(\Ideas, A)$. The chapter closes with a zoo of eight historical
positions on emergence --- Mill, Broad, Anderson, Bedau, Kim, Chalmers,
Holland, Sawyer --- each one a special case, refinement, or restriction of
the cocycle picture.

The mathematical content of the chapter has been formally verified in the
Lean module \texttt{IdeaTheory.Extensions.EmergenceLevels}; we cite specific
declarations from that module throughout. The sources of the proofs sketched
here are also collected, for the cumulative-chain identities, in the file
\texttt{Theorems5.lean} (the Structural Coherence and Hierarchical
Decomposition theorems) and, for the unit-insertion calculus, in
\texttt{Theorems7.lean} (the Composition Identity and Identity-Insertion
Stability theorems for transmission chains).

\section{The diagnostic: when does a valuation fail to be a homomorphism?}
\label{sec:diagnostic}

Fix a bare idea algebra $\IdMon$ and an abelian group $(A, +, 0)$. Often $A
= \mathbb{R}$, but the entire formal apparatus works for any abelian group,
and the case $A = \mathbb{Z}/2$ will reappear in
Chapter~\ref{ch:grading} when we discuss parity invariants.

\begin{definition}[Valuation]\label{def:valuation}
A \emph{valuation} on $\IdMon$ with values in $A$ is any function
\[
  v : \Ideas \longrightarrow A.
\]
A valuation is \emph{additive} (or a \emph{homomorphism}) when
\[
  v(\unit) = 0
  \qquad\text{and}\qquad
  v(a\compose b) \;=\; v(a) + v(b)
  \qquad\text{for all } a,b\in\Ideas.
\]
\end{definition}

Examples of intended valuations are: the description length of an idea, its
energy budget, its expected reward, its probability under a generative
model, its number of conceptual primitives, its reading time, its citation
count, the loss it incurs against a benchmark dataset, and so on. None of
these are additive in general. The description length of the conjunction of
two hypotheses is not the sum of their description lengths. The expected
reward of a composite policy is rarely the sum of the rewards of its
factors. The reading time of two arguments together exceeds the sum of
their separate reading times by an amount that depends on how much one
prepares for the other.

\begin{definition}[Emergence cocycle of a valuation]\label{def:omega}
For a valuation $v : \Ideas \to A$, the \emph{emergence cocycle} of $v$ is
the binary function $\omega_v : \Ideas\times\Ideas \to A$ defined by
\begin{equation}\label{eq:omega-def}
  \omega_v(a,b) \;:=\; v(a\compose b) \;-\; v(a) \;-\; v(b).
\end{equation}
\end{definition}

When the valuation is fixed by context we drop the subscript and write
$\omega(a,b)$. The Lean declaration is
\texttt{IdeaTheory.EmergenceLevels.omega}.

The cocycle is, by construction, the obstruction to additivity:

\begin{lemma}[Additivity criterion]\label{lem:additivity}
A valuation $v$ with $v(\unit) = 0$ is additive if and only if $\omega_v
\equiv 0$.
\end{lemma}

\begin{proof}
If $v$ is a homomorphism then $\omega_v(a,b) = v(a\compose b) - v(a) - v(b) =
0$ by definition; conversely, if $\omega_v \equiv 0$ then $v(a\compose b) =
v(a) + v(b)$, and $v(\unit) = 0$ holds by hypothesis.
\end{proof}

In Lean this is the \texttt{isAdditive\_of\_hom} declaration; the converse
direction is immediate from \texttt{IsAdditive} unfolding to the vanishing of
\texttt{omega}.

The diagnostic stance we take throughout the book is the following. Pick any
quantity of interest --- pick anything you might write down a number for ---
and ask: is its valuation additive over composition? Almost always the
answer is no, and almost always the failure is structured. The cocycle
$\omega_v$ is the structured failure. Far from a defect of the valuation, it
is a positive piece of data: \emph{$\omega_v$ tells us what new resources are
demanded, or new value is generated, by the act of composition.} A theory of
emergence is, on this view, simply a theory of $\omega$.

\begin{example}\label{ex:bool}
Let $\Ideas = \mathrm{Bool}$ with composition given by disjunction and unit
$\bot$. (The Lean instance for this idea algebra is in
\texttt{IdeaTheory.EmergenceLevels}.) Let $v(\bot) = 0$, $v(\top) = 1$. Then
\[
  \omega_v(\top,\top) \;=\; v(\top\vee\top) - v(\top) - v(\top) \;=\; 1 - 1
  - 1 \;=\; -1,
\]
so $v$ is not additive. We will return to this example throughout: it is the
minimal nontrivial idea algebra exhibiting strong emergence.
\end{example}

\section{The cocycle identity}\label{sec:cocycle}

The associativity of $\compose$ forces $\omega_v$ to satisfy a single
quadratic identity in three variables. This identity is the engine of every
subsequent result.

\begin{theorem}[Cocycle identity, $\partial\omega = 0$]\label{thm:cocycle}
Let $v : \Ideas \to A$ be any valuation and let $\omega = \omega_v$ be its
emergence cocycle. For all $a,b,c \in \Ideas$,
\begin{equation}\label{eq:cocycle}
  \omega(b,c) \;-\; \omega(a\compose b,\,c) \;+\; \omega(a,\,b\compose c)
  \;-\; \omega(a,b) \;=\; 0.
\end{equation}
\end{theorem}

\begin{proof}
We expand each term using \eqref{eq:omega-def}:
\begin{align*}
  \omega(b,c) &= v(b\compose c) - v(b) - v(c),\\
  \omega(a\compose b, c) &= v((a\compose b)\compose c) - v(a\compose b) -
  v(c),\\
  \omega(a, b\compose c) &= v(a\compose (b\compose c)) - v(a) - v(b\compose
  c),\\
  \omega(a,b) &= v(a\compose b) - v(a) - v(b).
\end{align*}
Substituting these into the left-hand side of \eqref{eq:cocycle} and
collecting like terms, the variables $v(a), v(b), v(c), v(a\compose b),
v(b\compose c)$ all cancel in matched $\pm$ pairs, and we are left with
\[
  v(a\compose (b\compose c)) \;-\; v((a\compose b)\compose c).
\]
By the associativity axiom for the monoid $\IdMon$, $(a\compose b)\compose c
= a\compose (b\compose c)$, so this difference is $v$ applied to the same
element minus itself, namely $0$.
\end{proof}

The cancellation is most easily verified by writing the eight $v$-terms in
two columns:
\[
\begin{array}{r|l}
\text{enters with sign }+ & \text{enters with sign }- \\\hline
v(b\compose c),\ v(a\compose((b\compose c))),\ v(a\compose b),\ v(a),\ v(b)
& v(b),\ v(c),\ v((a\compose b)\compose c),\ v(a\compose b),\ v(c),\\
& v(a),\ v(b\compose c),\ v(a\compose b),\ v(a),\ v(b)
\end{array}
\]
After cancelling the duplicated terms one is left with the difference
$v(a\compose(b\compose c)) - v((a\compose b)\compose c)$, which vanishes by
associativity. The verified declaration is
\texttt{IdeaTheory.EmergenceLevels.cocycle\_identity}.

\begin{remark}[The name ``cocycle'']\label{rem:cocycle}
Write $C^n(\Ideas, A)$ for the abelian group of functions $\Ideas^n \to A$,
and define the coboundary maps $\partial^n : C^n \to C^{n+1}$ by the usual
bar-resolution formulas:
\begin{align*}
  (\partial^0 a)(x) &= 0 \qquad (a\in A,\ x\in\Ideas),\\
  (\partial^1 f)(a,b) &= f(a\compose b) - f(a) - f(b),\\
  (\partial^2 \omega)(a,b,c) &= \omega(b,c) - \omega(a\compose b, c) +
  \omega(a, b\compose c) - \omega(a,b).
\end{align*}
Theorem~\ref{thm:cocycle} is the assertion that $\partial^2 \omega_v = 0$ for
every valuation $v$, equivalently that the natural map $v \mapsto \omega_v$
factors through the kernel of $\partial^2$. This places $\omega_v$ as a
2-cocycle in the bar complex of the monoid $\IdMon$ with coefficients in
$A$, in the sense of monoid cohomology. We will not need the full machinery
of $H^*(\IdMon, A)$, but we will use $H^2$ in Section~\ref{sec:bounded}.
\end{remark}

\begin{remark}[Why associativity is the right hypothesis]
The proof of Theorem~\ref{thm:cocycle} uses \emph{exactly} associativity ---
nothing else. It uses neither the existence of a unit, nor the algebra
operations on $A$ beyond addition, nor any property of $v$ beyond its being
a function. The cocycle identity is therefore a fact about associative
binary operations and abelian groups, not a fact peculiar to ideas. This is
why every quantitative theory whose underlying combinatorics is associative
ends up, eventually, encountering a 2-cocycle: physics with central
extensions of Lie algebras, projective representations of groups, the
Schur multiplier, the Brauer group, the Hochschild $H^2$ of an algebra. The
emergence cocycle of an idea algebra is a member of this family.
\end{remark}

\section{Normalization}\label{sec:normalization}

The cocycle of a generic valuation $v$ has nonzero ``edge values'' coming
from $v(\unit)$. We record the precise form of these.

\begin{lemma}[Edge values of $\omega$]\label{lem:edge}
For any valuation $v : \Ideas \to A$ and any $a \in \Ideas$,
\[
  \omega_v(\unit, a) \;=\; -\,v(\unit), \qquad
  \omega_v(a, \unit) \;=\; -\,v(\unit).
\]
\end{lemma}

\begin{proof}
By the unit law, $\unit\compose a = a$, so $\omega_v(\unit, a) = v(\unit
\compose a) - v(\unit) - v(a) = v(a) - v(\unit) - v(a) = -v(\unit)$. The
right-edge calculation is symmetric, using $a\compose \unit = a$.
\end{proof}

Hence:

\begin{definition}[Normalized valuation]
A valuation $v$ is \emph{normalized} if $v(\unit) = 0$.
\end{definition}

\begin{corollary}[Normalization of $\omega$]\label{cor:normalized-omega}
If $v$ is normalized then
\[
  \omega_v(\unit, a) = 0 \qquad\text{and}\qquad \omega_v(a,\unit) = 0
  \qquad\text{for all } a\in\Ideas.
\]
\end{corollary}

In Lean these are \texttt{omega\_ident\_left\_of\_zero} and
\texttt{omega\_ident\_right\_of\_zero}. Together they say that the cocycle
of a normalized valuation is a \emph{normalized 2-cocycle} in the sense of
the standard normalized bar complex of $\IdMon$.

Normalization is mild: every valuation $v$ can be normalized by passing to
$\widetilde v(a) := v(a) - v(\unit)\cdot\chi(a)$, where $\chi$ is any
1-cochain with $\chi(\unit) = 1$. We will discuss this rigidification more
carefully in Section~\ref{sec:coboundary}. For the rest of the chapter we
silently assume valuations are normalized; the reader who prefers not to may
restore the constant $-v(\unit)$ at the unit edges by hand.

\begin{remark}
The choice of normalization corresponds, in the bar complex, to working with
the reduced (or normalized) chain groups $\bar C^n$, where $n$-cochains
vanish whenever any argument is the unit. The inclusion $\bar C^* \hookrightarrow
C^*$ is a quasi-isomorphism, so $H^*(\IdMon, A)$ may be computed either
way.
\end{remark}

\section{Coboundaries and the equivalence relation on cocycles}
\label{sec:coboundary}

Two distinct valuations may carry the same emergence content, in the
following precise sense.

\begin{definition}[Coboundary of a 1-cochain]\label{def:coboundary}
For $f : \Ideas \to A$, the \emph{coboundary} of $f$ is the function
$\delta f : \Ideas\times\Ideas \to A$ defined by
\[
  (\delta f)(a,b) \;:=\; f(a\compose b) - f(a) - f(b).
\]
A function $\omega : \Ideas\times\Ideas \to A$ is a \emph{coboundary} if
$\omega = \delta f$ for some $f$.
\end{definition}

The Lean definition is \texttt{IdeaTheory.EmergenceLevels.delta}; observe
that, by Definition~\ref{def:omega}, $\omega_v = \delta v$. So:

\begin{lemma}\label{lem:cob-of-val}
The cocycle of a valuation $v$ is the coboundary of $v$ regarded as a
1-cochain. Symbolically, $\omega_v = \delta v$.
\end{lemma}

In particular, every $\omega_v$ is a coboundary in the bar complex of
$\IdMon$, hence is cohomologically trivial. This is not a contradiction
with Section~\ref{sec:bounded}, where we will reintroduce nontrivial
cohomology classes via a more refined notion of valuation. The point is:
\emph{within the class of valuations $v : \Ideas \to A$ with no further
structure}, every cocycle is a coboundary --- because every cocycle
\emph{came from} a 1-cochain by Definition~\ref{def:omega}. The relevant
question is which $\omega$ \emph{cannot be made to vanish} by replacing $v$
with another valuation in a constrained class. To formalize this we look at
how $\omega_v$ shifts when we perturb $v$.

\begin{proposition}[Shift formula]\label{prop:cob}
For any valuations $v, w : \Ideas \to A$,
\begin{equation}\label{eq:shift}
  \omega_{v+w}(a,b) \;=\; \omega_v(a,b) + \omega_w(a,b).
\end{equation}
In particular, if $w = f$ is regarded as a 1-cochain, then $\omega_{v+f} =
\omega_v + \delta f$.
\end{proposition}

\begin{proof}
Direct expansion:
\begin{align*}
  \omega_{v+w}(a,b) &= (v+w)(a\compose b) - (v+w)(a) - (v+w)(b)\\
  &= [v(a\compose b) - v(a) - v(b)] + [w(a\compose b) - w(a) - w(b)]\\
  &= \omega_v(a,b) + \omega_w(a,b). \qedhere
\end{align*}
\end{proof}

The Lean form is \texttt{omega\_add} together with \texttt{omega\_sum}. The
moral is that the map $v \mapsto \omega_v$ is $A$-linear in $v$, and it is
exactly the coboundary operator $\delta$ of the bar complex.

\begin{definition}[Cohomologous cocycles]\label{def:cohomologous}
Two cocycles $\omega, \omega' \in \ker\partial^2$ are \emph{cohomologous},
written $\omega \sim \omega'$, if their difference is a coboundary:
\[
  \omega \sim \omega' \quad\Longleftrightarrow\quad \omega' - \omega = \delta
  f \text{ for some } f : \Ideas\to A.
\]
The set of equivalence classes is
\[
  H^2(\IdMon, A) \;:=\; \ker\partial^2 \,/\, \operatorname{im}\partial^1.
\]
\end{definition}

\begin{remark}[The cohomology class of a valuation cocycle]\label{rem:class-trivial}
Lemma~\ref{lem:cob-of-val} says that the class $[\omega_v] \in H^2(\IdMon, A)$
of any valuation-induced cocycle is trivial. This sounds disappointing.
The way nontriviality re-enters is by enlarging the universe of cocycles:
we admit any $\omega : \Ideas^2 \to A$ that satisfies $\partial^2 \omega = 0$
and is normalized, regardless of whether it arises from some $v$. In
Sections~\ref{sec:bounded} and \ref{sec:zoo} we will see that such free
cocycles correspond exactly to the strong-emergence phenomena of Chalmers
\cite{ChalmersStrong} and the irreducible group cognition of Sawyer
\cite{Sawyer2005}: there exists structure on $\Ideas\times\Ideas$ that no
1-dimensional valuation can capture.
\end{remark}

\section{Telescoping along chains}\label{sec:telescope}

Composition of three or more ideas can be parenthesized in many ways; by
associativity all parenthesizations agree as elements of $\Ideas$, but their
\emph{cocycle accountings} look different. The Transmission Telescoping
Theorem reconciles all of them into a single formula.

\begin{notation}\label{not:chain}
For a list $a_1, \dots, a_n \in \Ideas$, write
\[
  a_1 \compose a_2 \compose \cdots \compose a_n
\]
for the (unambiguous, by associativity) composition of all $n$ factors.
The Lean implementation is the \texttt{CompChain} function from
\texttt{Theorems5.lean}; the basic computational identities (e.g.\ that
\texttt{CompChain [a,b]} equals \texttt{op a b}) are
\texttt{comp\_chain\_two}, \texttt{comp\_chain\_three}, etc.
\end{notation}

We require one piece of book-keeping: the \emph{prefix products}
\[
  P_k \;:=\; a_1 \compose a_2 \compose \cdots \compose a_k, \qquad 1\le k\le n,
\]
with the convention $P_1 = a_1$.

\begin{theorem}[Transmission Telescoping]\label{thm:telescope}
Let $v : \Ideas \to A$ be a normalized valuation with cocycle $\omega =
\omega_v$, and let $a_1,\dots,a_n \in \Ideas$ be any list with $n \ge 1$.
Then
\begin{equation}\label{eq:telescope}
  v\bigl(a_1\compose\cdots\compose a_n\bigr) \;=\; \sum_{k=1}^n v(a_k)
  \;+\; \sum_{k=1}^{n-1} \omega(P_k,\,a_{k+1}).
\end{equation}
\end{theorem}

In words: \emph{the value of a chained composition is the sum of the values
of the parts plus a sum of pairwise emergence terms, one for each
junction.} The emergence terms are computed against the running prefix, so
the cocycle is sensitive to the entire history at the point of each new
junction.

\begin{proof}
We argue by induction on $n$. The base cases $n=1$ and $n=2$ are immediate:
for $n=1$, both sides equal $v(a_1)$ and the cocycle sum is empty; for
$n=2$, $P_1 = a_1$ and \eqref{eq:telescope} reads $v(a_1\compose a_2) =
v(a_1) + v(a_2) + \omega(a_1, a_2)$, which is the definition of $\omega$.

For the inductive step, suppose \eqref{eq:telescope} holds for chains of
length $n$. Let $a_1, \dots, a_{n+1}$ be a chain of length $n+1$, and write
$Q := P_n = a_1\compose\cdots\compose a_n$, so that the full product is
$Q\compose a_{n+1}$. By the cocycle definition applied at $(Q, a_{n+1})$,
\[
  v(Q\compose a_{n+1}) \;=\; v(Q) + v(a_{n+1}) + \omega(Q, a_{n+1}).
\]
By the inductive hypothesis applied to $Q$,
\[
  v(Q) \;=\; \sum_{k=1}^n v(a_k) + \sum_{k=1}^{n-1} \omega(P_k, a_{k+1}).
\]
Substituting,
\[
  v(Q\compose a_{n+1}) \;=\; \sum_{k=1}^n v(a_k) + \sum_{k=1}^{n-1}
  \omega(P_k, a_{k+1}) + v(a_{n+1}) + \omega(P_n, a_{n+1}).
\]
Reindexing the sums to combine $v(a_{n+1})$ into the first sum and the new
$\omega$-term into the second sum, this is precisely
\[
  \sum_{k=1}^{n+1} v(a_k) + \sum_{k=1}^{n} \omega(P_k, a_{k+1}),
\]
which is \eqref{eq:telescope} for the length-$(n+1)$ chain.
\end{proof}

\begin{corollary}[Symmetric form]\label{cor:tele-sym}
If $\omega$ is, in addition, \emph{symmetric}, i.e.\ $\omega(a,b) =
\omega(b,a)$, then equation \eqref{eq:telescope} can be rewritten as
\[
  v(a_1\compose\cdots\compose a_n) \;=\; \sum_{k=1}^n v(a_k) + \sum_{1\le i
  < j \le n} \omega(a_i, a_j),
\]
provided $\omega$ extends additively in each argument over composition (an
extra hypothesis that holds, e.g., when $\omega(a,b)$ depends only on
``content tags'' of $a,b$ that are themselves additive).
\end{corollary}

\begin{remark}[Choice of parenthesization]
The statement of Theorem~\ref{thm:telescope} appears to depend on the
left-associative parenthesization $((a_1\compose a_2)\compose
a_3)\cdots$. It does not. By Theorem~\ref{thm:cocycle} and the universality
of associativity, any other parenthesization yields the same total emergence
sum, modulo a permutation of the cocycle terms that the cocycle identity
\eqref{eq:cocycle} forces to cancel. A formal proof uses the
\texttt{decomp\_unique\_structure} lemma of \texttt{Theorems5.lean}: any two
chain decompositions of a fixed element of $\Ideas$ have equal
\texttt{CompChain} value, and hence equal $v$-image, so the right-hand sides
must agree.
\end{remark}

The telescoping theorem is the precise version of the informal claim, made
in nearly every emergentist text since Mill, that ``compound effects build
up over a sequence of combinations.'' We will see in
Sections~\ref{sec:zoo:holland} and \ref{sec:zoo:henrich-style} that
several seemingly disparate models of cumulative novelty --- Holland's
adaptive agents \cite{Holland1998}, Henrich's cultural ratchet
\cite{Henrich2015} --- are precisely \eqref{eq:telescope} with different
choices of $A$ and $\omega$.

\subsection*{Functoriality and a glance at higher coboundaries}

The shift formula \eqref{eq:shift} has a categorical reading worth
recording. The assignment $A \mapsto C^*(\IdMon, A)$ from abelian groups to
cochain complexes is functorial: a homomorphism $\varphi : A \to A'$
induces $\varphi_* : C^n(\IdMon, A) \to C^n(\IdMon, A')$ by
post-composition, and $\varphi_*$ commutes with $\partial$. Consequently,
the cocycle of a transported valuation is the transport of the cocycle:
\[
  \omega_{\varphi\circ v} \;=\; \varphi \circ \omega_v.
\]
We will use this in Chapter~\ref{ch:grading} when transporting valuations
along the grading homomorphism $\deg : \Ideas \to G$ and along quotient
maps $G \twoheadrightarrow G/H$ that ``forget'' parts of the degree.

A second piece of structural overhead, used only briefly in the present
chapter, is the next coboundary $\partial^3 : C^3 \to C^4$:
\[
  (\partial^3 \xi)(a,b,c,d) := \xi(b,c,d) - \xi(a\compose b, c, d) +
  \xi(a, b\compose c, d) - \xi(a, b, c\compose d) + \xi(a,b,c).
\]
A short calculation, of the same flavor as the proof of
Theorem~\ref{thm:cocycle}, gives $\partial^3 \circ \partial^2 = 0$, so the
$C^*(\IdMon, A)$ form a genuine cochain complex. The class $[\omega_v] \in
H^2(\IdMon, A)$ is therefore a well-defined invariant attached to the
isomorphism class of $v$ in the category of valuations modulo
1-coboundaries.

\subsection*{Three worked examples of telescoping}\label{sec:tele-examples}

We illustrate Theorem~\ref{thm:telescope} on three idea algebras of
increasing realism.

\begin{example}[Powerset, cardinality]\label{ex:tele-power}
Let $\Ideas = \mathbf{P}(X)$ for a finite set $X$, with $\compose = \cup$,
$\unit = \emptyset$, and $v(S) = |S|$. Then $\omega_v(S, T) = -|S\cap T|$
(Example~\ref{ex:signs}). For sets $S_1, \dots, S_n$, the prefix products
are $P_k = S_1 \cup \cdots \cup S_k$ and \eqref{eq:telescope} reads
\[
  \Bigl|\bigcup_{k=1}^n S_k\Bigr| \;=\; \sum_{k=1}^n |S_k| \;-\;
  \sum_{k=1}^{n-1} |P_k \cap S_{k+1}|,
\]
which is the standard step-wise inclusion-exclusion formula. The cocycle is
sub-additive throughout, so by Lemma~\ref{lem:sign-tele} the cardinality
of a union is at most the sum of the cardinalities, with equality iff the
$S_k$ are pairwise disjoint.
\end{example}

\begin{example}[Free monoid, length]\label{ex:tele-free}
Let $\Ideas = \Sigma^*$ be the free monoid on an alphabet $\Sigma$, with
$\compose$ concatenation and $\unit$ the empty string. Take $v(w) =
\mathrm{length}(w)$. Then $v$ is additive, $\omega_v \equiv 0$, and the
telescoping formula collapses to the trivial $|w_1\cdots w_n| = \sum_k
|w_k|$. So the free monoid carries no length-emergence: it is the formal
home of a fully-reductive valuation.
\end{example}

\begin{example}[Free monoid, Kolmogorov complexity]\label{ex:tele-K}
Same $\Ideas = \Sigma^*$, but now let $v(w) = K(w)$ be a fixed prefix-free
Kolmogorov complexity. The cocycle $\omega_K(u, v) = K(uv) - K(u) - K(v)$ is
no longer zero; up to an $O(\log\min(K(u), K(v)))$ term it is the negative
of the algorithmic mutual information $I(u : v)$, by the symmetry-of-
information theorem. The telescoping formula here gives
\[
  K(w_1\cdots w_n) \;=\; \sum_k K(w_k) \;-\; \sum_{k=1}^{n-1}
  I(P_k : w_{k+1}) + O(\log n),
\]
which is a recognizable identity in algorithmic information theory and a
direct application of Theorem~\ref{thm:telescope}.
\end{example}

These three examples exhibit, respectively, sub-additive bounded emergence,
no emergence, and the typical regime of unbounded but
information-theoretically-controlled emergence.

\section{Bounded vs.\ unbounded $\omega$: weak and strong emergence}
\label{sec:bounded}

For the rest of the chapter take $A = \mathbb{R}$. The two modes of
emergence universally distinguished in the philosophical literature ---
weak and strong --- correspond to two quantitative properties of $\omega$.

\begin{definition}[Weak/bounded emergence]\label{def:weak}
A valuation $v : \Ideas \to \mathbb{R}$ is \emph{weakly emergent} (or
\emph{$\omega$-bounded}) with bound $M \ge 0$ if
\[
  |\omega_v(a,b)| \;\le\; M \qquad\text{for all } a, b \in \Ideas.
\]
The infimum of admissible $M$ is the \emph{emergence radius} of $v$, written
$\|\omega_v\|_\infty$.
\end{definition}

The Lean predicate is \texttt{IsBoundedEmergent}; the lemmas
\texttt{isBoundedEmergent\_zero\_of\_additive},
\texttt{isBoundedEmergent\_mono}, \texttt{isBoundedEmergent\_nonneg},
\texttt{isBoundedEmergent\_smul}, \texttt{isBoundedEmergent\_add},
\texttt{isBoundedEmergent\_add\_self} establish, respectively, that
additive valuations are weakly emergent with $M=0$; the bound is monotone;
the bound is non-negative on nonempty algebras; non-negative scalar
multiples scale the bound; and the bound is sub-additive under addition of
valuations.

\begin{definition}[Strong emergence]\label{def:strong}
A valuation $v$ is \emph{strongly emergent} if there is no additive
valuation $v' : \Ideas \to \mathbb{R}$ that agrees with $v$ pointwise:
\[
  \neg\,\exists\,v'\,:\,\Ideas\to\mathbb{R} \;\text{ such that }\;
  v'\text{ is additive and } v'(a) = v(a) \text{ for all } a\in\Ideas.
\]
\end{definition}

The Lean predicate is \texttt{IsStronglyEmergent}.

The two definitions appear unrelated --- one quantitative, one
existential. The following observation, verified in Lean as
\texttt{stronglyEmergent\_iff\_not\_additive}, collapses the existential.

\begin{proposition}[Strong emergence as non-additivity]\label{prop:strong-iff}
A valuation $v$ is strongly emergent if and only if $v$ itself is not
additive, i.e.\ $\omega_v \not\equiv 0$.
\end{proposition}

\begin{proof}
The reverse direction is contrapositive: if $v$ is additive, then $v' := v$
witnesses the negation of strong emergence.

For the forward direction, suppose $v$ is not strongly emergent. Then there
is an additive $v'$ with $v'(a) = v(a)$ for all $a$. From pointwise
agreement, $v'(a\compose b) = v(a\compose b)$, $v'(a) = v(a)$, $v'(b) =
v(b)$. Substituting into $\omega_{v'}(a,b) = 0$ gives $v(a\compose b) -
v(a) - v(b) = 0$, i.e.\ $\omega_v(a,b) = 0$. Since $a,b$ were arbitrary,
$v$ is additive.
\end{proof}

This proposition is at first sight a deflation: ``strong emergence'' as a
predicate of \emph{single} valuations reduces to non-additivity, which is
the same as $\omega_v \neq 0$. The genuine cohomological content of strong
emergence appears only when one passes to the \emph{class of all valuations
sharing some structural property} --- e.g.\ ``valuations measurable with
respect to a sub-$\sigma$-algebra of constituents,'' or ``valuations
expressible as polynomial functions of base-feature counts,'' or, in the
philosophical setting, ``valuations expressible by physical
science''. Within such a constrained class, an $\omega$ that cannot be
realized as $\delta f$ for any $f$ \emph{in the class} is a genuine
cohomological obstruction. This is the precise reading we propose for
Chalmers \cite{ChalmersStrong}: weak emergence is bounded $\omega_v$ for an
allowed valuation $v$; strong emergence is an $\omega$-class in the
restricted cohomology $H^2(\IdMon, \mathbb{R})_{\mathrm{phys}}$ that is
nontrivial even though it may be a coboundary in the unrestricted
$H^2(\IdMon, \mathbb{R})$.

\begin{example}[Continuation of Example~\ref{ex:bool}]\label{ex:bool-strong}
The \texttt{Bool} indicator valuation $v(\bot) = 0$, $v(\top) = 1$ on the
disjunction monoid is weakly emergent with optimal bound $M = 1$ (verified
as \texttt{boolIndicator\_boundedEmergent}) and strongly emergent (verified
as \texttt{boolIndicator\_stronglyEmergent}). The key cocycle value is
$\omega_v(\top, \top) = -1$, witnessing non-additivity. So a single
two-element idea algebra suffices to realize both emergence modes. The
generalization to $\mathbf{P}(X)$ (the powerset of a finite set, with
union as composition and the cardinality $v(S) = |S|$ as valuation) is the
inclusion-exclusion principle, with $\omega_v(S, T) = -|S\cap T|$.
\end{example}

\begin{remark}[The taxonomy at a glance]
We may summarize the four positions a valuation can occupy as follows.
\begin{center}
\begin{tabular}{lll}\toprule
& $\omega_v \equiv 0$ & $\omega_v \not\equiv 0$\\\midrule
$|\omega_v| \le M$ for some $M$ & additive (no emergence) & weakly emergent\\
$|\omega_v|$ unbounded & impossible & strongly emergent (and unbounded)\\
\bottomrule
\end{tabular}
\end{center}
The upper-left cell is the trivial case; the upper-right is Bedau's regime;
the lower-right is Chalmers's regime; the lower-left is empty because
$\omega_v \equiv 0$ already gives bound zero. There is also a refinement of
the upper-right by the size of $\|\omega_v\|_\infty$: a valuation with very
small emergence radius is ``almost additive'' and we will sometimes call
such valuations \emph{$\varepsilon$-additive}.
\end{remark}

\subsection*{Sub-additivity, super-additivity, and the sign of $\omega$}

In many applied settings the cocycle has a definite sign. We pause to record
the standard vocabulary, which appears throughout the zoo of
Section~\ref{sec:zoo}.

\begin{definition}[Sign regimes]\label{def:sign}
A valuation $v$ is \emph{super-additive} if $\omega_v(a,b) \ge 0$ for all
$a, b \in \Ideas$, and \emph{sub-additive} if $\omega_v(a,b) \le 0$ for
all $a, b$. It is \emph{$\omega$-positive on $S \subseteq \Ideas\times
\Ideas$} if $\omega_v|_S \ge 0$, and analogously for negative.
\end{definition}

\begin{example}[Common sign patterns]\label{ex:signs}
\textit{(a)} Cardinality on the powerset monoid $(\mathbf{P}(X), \cup,
\emptyset)$ is sub-additive: $\omega_v(S, T) = -|S \cap T| \le 0$, with
equality iff $S$ and $T$ are disjoint. \textit{(b)} The free-energy of a
joint physical system is super-additive when interactions are repulsive and
sub-additive when interactions are attractive. \textit{(c)} Description
length under any prefix-free code is super-additive on average and
sub-additive in the presence of shared sub-strings, by Kraft's inequality.
\textit{(d)} The Boolean indicator of Example~\ref{ex:bool} is sub-additive
with $\omega_v(\top,\top) = -1$.
\end{example}

\begin{lemma}[Sign and telescoping]\label{lem:sign-tele}
If $v$ is super-additive (resp.\ sub-additive), then for every chain
$a_1, \dots, a_n$ with $n \ge 1$,
\[
  v(a_1\compose\cdots\compose a_n) \;\ge\; \sum_{k=1}^n v(a_k)
  \quad(\text{resp. } \le).
\]
\end{lemma}

\begin{proof}
By Theorem~\ref{thm:telescope} the left-hand side equals $\sum_k v(a_k) +
\sum_{k=1}^{n-1} \omega_v(P_k, a_{k+1})$. If every $\omega_v(P_k, a_{k+1})
\ge 0$ then the cocycle sum is non-negative, and the inequality follows;
the sub-additive case is dual.
\end{proof}

The sign regime is preserved under non-negative scaling and under sums of
same-sign valuations. It interacts well with bounded emergence: a
super-additive $M$-bounded valuation has cocycle in $[0, M]$, which is the
appropriate setting for entropy-like and complexity-like valuations. The
sign vocabulary will be most useful in Section~\ref{sec:zoo:anderson} (where
positive cocycle is what makes ``more'' different), in
Section~\ref{sec:zoo:holland} (where same-sign cocycles cumulate
monotonically along chains), and in Section~\ref{sec:zoo:mill} (where
heteropathy is silent on sign and so sub-additive and super-additive
heteropathy are both observed empirically).

\section{The zoo}\label{sec:zoo}

We close with eight historical positions on emergence, each one a special
case, refinement, or restriction of the cocycle picture.

\subsection{Mill --- heteropathic laws}\label{sec:zoo:mill}

\paragraph{Original claim.} In Book III, Chapter VI of \emph{A System of
Logic} \cite{Mill1843}, John Stuart Mill distinguishes
\emph{homopathic} from \emph{heteropathic} laws of causation. A homopathic
law is one in which ``the joint effect of several causes is identical with
the sum of their separate effects''; this is the principle Mill calls the
\emph{Composition of Causes}, and it is exemplified by the parallelogram of
forces in mechanics. A heteropathic law is one in which the principle
fails: the conjoint operation of two causes produces an effect that is not
the resultant of what each would have produced alone. Mill cites chemistry
as the paradigmatic source of heteropathic laws --- the properties of
water are not the sum of the properties of hydrogen and oxygen --- and
extends the morphology to physiological and mental phenomena. The
distinction is meant to mark the boundary between the domain of
deductive-mechanical science and the domain of what later authors will
call emergent science.

\paragraph{Formal restatement.} Take $\Ideas$ to be the set of causes (or
mechanisms) under study, $\compose$ to be their conjoint operation, and let
$v : \Ideas \to A$ be the valuation that records ``the effect''. Mill's
\emph{Composition of Causes} is exactly the assertion that $v$ is additive,
i.e.\ $\omega_v \equiv 0$. Heteropathy is the assertion that $\omega_v$ is
not identically zero somewhere, i.e.\ that there exist $a, b \in \Ideas$
with $\omega_v(a, b) \ne 0$. Mill's chemical example is the chemical
analogue of Example~\ref{ex:bool-strong}: at the pair (hydrogen, oxygen)
the cocycle takes a nonzero value (the value, in fact, that yields water).

\paragraph{Status.} Made true (\emph{ipso facto}, by definition) by
Theorem~\ref{thm:cocycle} together with Lemma~\ref{lem:additivity}: every
valuation has a cocycle, the cocycle is identically zero exactly when the
valuation is homopathic, and otherwise the cocycle records, pointwise, the
extent and direction of heteropathy. So Mill's typology is not a typology of
laws but of \emph{(valuation, pair-of-causes)} relations, and
$\omega_v(a,b)$ is the precise quantitative content of ``how much
heteropathy'' obtains at $(a,b)$.

\paragraph{Commentary.} Mill himself worried that heteropathic laws might
admit no scientific treatment, since the joint effect could be ``anything
whatever''. The cocycle picture corrects this by exhibiting the
\emph{constraint} on heteropathy: even maximally heteropathic $\omega_v$
must satisfy the cocycle identity $\partial^2\omega_v = 0$
\eqref{eq:cocycle}, so heteropathic laws are not lawless --- they are
laws of a different (cohomological) type. Mill's insight is correct in
direction but understates the residual structure available. We will use this
correction in Chapter~\ref{ch:grading} to show that even thoroughly
heteropathic phenomena may admit a coherent grading.

\subsection{C.\,D.\ Broad --- emergent qualities}\label{sec:zoo:broad}

\paragraph{Original claim.} In \emph{The Mind and Its Place in Nature}
\cite{Broad1925}, C.\,D.\ Broad articulates the emergentist position with
greater precision than Mill. Broad's central concept is that of an
\emph{emergent quality}: a property of an aggregate $A_1 + \cdots + A_n$
that is not deducible, even by an ideal scientist, from the most complete
knowledge of the properties of the $A_i$ in isolation \emph{and} of the
properties of all aggregates of comparable kind in which the $A_i$
appear. Broad distinguishes mechanism (everything is in principle
deducible), emergent materialism (some properties are emergent in his
strict sense), and substantival vitalism (the aggregate has constituents not
present in the parts). His thesis, applied to chemistry, biology, and mind,
is the second.

\paragraph{Formal restatement.} Take $A = \mathbb{R}$ and consider a
restricted class $\mathcal V_{\mathrm{deduc}}$ of valuations $v$ that are
``deducible from constituents'' --- formally, those expressible as the
restriction to $\Ideas$ of some prescribed set of formulas in the
properties of atomic ideas. Broad's claim is that there are pairs $(a, b)$
and quantities $v$ such that $\omega_v(a, b) \ne 0$ \emph{and} $v$ does
not lie in $\mathcal V_{\mathrm{deduc}}$. Equivalently, $[\omega_v] \in
H^2(\IdMon, \mathbb{R})_{\mathcal V_{\mathrm{deduc}}}$ is nontrivial.

\paragraph{Status.} Contingent on the choice of $\mathcal
V_{\mathrm{deduc}}$. The cocycle picture establishes that the question
``which valuations have nontrivial emergence?'' is a question about
$H^2$ relative to a class of allowed cobounding 1-cochains. The Lean
predicate \texttt{IsStronglyEmergent} witnesses Broad's notion at the
\emph{absolute} level (no additive $v'$ at all matches $v$ pointwise);
Broad's interest is the relative case. Theorem~\ref{thm:cocycle} forces
even ``emergent qualities'' to obey the cocycle identity, which Broad did
not anticipate.

\paragraph{Commentary.} Broad's careful definition of ``deducibility'' as
relative to a base of \emph{accessible aggregates} is the precise content of
the constraint that the cobounding 1-cochain $f$ live in a particular
function class. As McLaughlin and others have noted in the revival of
British emergentism \cite{Anderson1972}, the trouble with classical
emergentism was the absence of a vocabulary for the constraint. The cocycle
picture supplies it: Broad's emergent qualities are cocycles whose
classes survive in the obstruction group $H^2$ of the deducibility
calculus.

\subsection{Anderson --- ``More is different''}\label{sec:zoo:anderson}

\paragraph{Original claim.} Philip W.\ Anderson's 1972 \emph{Science}
article \cite{Anderson1972} argues against the constructionist hypothesis ---
the view that knowledge of fundamental laws plus computational power
suffices to derive the behavior of large aggregates. Anderson's slogan,
``more is different,'' asserts that at each level of complexity new
phenomena appear that require new concepts. The article is the founding
manifesto of condensed-matter and complex-systems perspectives in physics,
and supplies the rallying cry of the emergentist research program for the
remainder of the twentieth century.

\paragraph{Formal restatement.} Anderson is asserting that the iterated
composition $a_1 \compose \cdots \compose a_n$ admits valuations $v$ for
which the residual sum $\sum_{k=1}^{n-1} \omega_v(P_k, a_{k+1})$ from
Theorem~\ref{thm:telescope} dominates the linear term $\sum_{k=1}^n
v(a_k)$ as $n$ grows. ``More is different'' is the assertion that
$\omega_v$ is not only nonzero but unbounded, so that the cocycle
contribution to large-$n$ chains becomes the leading-order phenomenon. In
this formulation the magnitude of $v(a_1\compose\cdots\compose a_n)$ is
not predicted by knowledge of $v(a_k)$ alone.

\paragraph{Status.} True under the explicit hypothesis that $\omega_v$ is
unbounded; false otherwise. Theorem~\ref{thm:telescope} makes the slogan
\emph{quantitative}: it says exactly how much more is different, namely the
prefix-sum $\sum \omega_v(P_k, a_{k+1})$. The Lean lemma
\texttt{isBoundedEmergent\_nonneg} forces this sum to be at most $(n-1)M$
when $\omega_v$ is $M$-bounded; in the strong-emergence regime the bound
fails and the sum can grow super-linearly in $n$.

\paragraph{Commentary.} Anderson's article is often misread as a
\emph{philosophical} thesis about reduction. Read against the cocycle
picture, it is in fact a \emph{quantitative} claim about a specific limit:
in the regime $n \to \infty$, the cocycle term in \eqref{eq:telescope}
overwhelms the additive term. The cocycle picture predicts when this
happens (whenever $|\omega_v|$ has a positive lower bound on a sufficiently
recurrent sub-monoid) and when it does not (whenever $\omega_v$ is summable
in some norm). This makes ``more is different'' a testable hypothesis per
choice of $\Ideas$ and $v$, not a metaphysical posit.

\subsection{Bedau --- weak emergence}\label{sec:zoo:bedau}

\paragraph{Original claim.} Mark Bedau \cite{Bedau1997} introduces the
concept of \emph{weak emergence}: a macrostate $P$ of a system $S$ is
weakly emergent just in case $P$ is \emph{derivable} from $S$'s
microdynamics, but only by simulation --- i.e., there is no shortcut
deduction, only step-by-step iteration. Bedau's notion is intended to
respect the autonomy of the special sciences (which study weakly emergent
properties) without demanding the metaphysical commitments of strong
emergentism. Subsequent literature \cite{Holland1998} adopts weak
emergence as the standard formal notion in complex-systems modeling.

\paragraph{Formal restatement.} Bedau's regime is the case where $\omega_v$
exists, is nonzero, and is uniformly bounded: there is some $M$ with
$|\omega_v(a,b)| \le M$ for all pairs. By Theorem~\ref{thm:telescope},
the value of a chain of length $n$ is then known up to additive error at
most $(n-1)M$ once one knows the parts $a_k$. ``Derivable by simulation''
becomes ``compute the prefix products $P_k$ and sum the cocycle terms ---
no closed-form is possible without doing the computation.''

\paragraph{Status.} True \emph{by definition} once the verbal definition is
translated into the cocycle vocabulary. The Lean predicate
\texttt{IsBoundedEmergent} is a literal mathematization of Bedau's
informal definition, with the bound $M$ playing the role of the ``small
constant of unpredictability per simulation step.''

\paragraph{Commentary.} Bedau's appeal to ``simulation'' has a precise
algebraic shadow in the cocycle picture: the simulation is the iterated
evaluation of $\omega_v$ at successive prefix-junctions $P_k$, and the
``inevitability of simulation'' is the algebraic fact that $\omega_v$ does
not factor through a hom (i.e., is not a coboundary of an
``algorithmically-evaluable'' 1-cochain $f$). The weak-emergence regime is
also the regime in which the cocycle contributes only a bounded correction
per step, so that ergodic averages of $v$ along chains are well-defined.
This is the technical content of the ``coarse-graining'' perspectives in
complex-systems modeling, which are recovered as quotients of $\Ideas$
under equivalence by chains whose cumulative cocycle is bounded.

\subsection{Kim --- downward causation and the exclusion problem}
\label{sec:zoo:kim}

\paragraph{Original claim.} Jaegwon Kim's \emph{Mind in a Physical World}
\cite{Kim1998} formulates the \emph{exclusion problem} for nonreductive
physicalism. The problem is this: if every physical effect has a sufficient
physical cause, and if mental properties supervene on physical ones, then
mental causes are causally redundant unless mental properties have causal
powers \emph{not} possessed by their physical realizers. The latter would be
\emph{downward causation}, and Kim argues that downward causation is
incoherent under standard supervenience. The result is a dilemma: either
mental properties are reducible (and the special sciences have no
autonomous causal vocabulary), or they are causally inert.

\paragraph{Formal restatement.} Take $\Ideas$ to be a category of physical
states with composition the temporal evolution operator, and $v$ a
``mental'' valuation. Supervenience is the assertion that $v$ factors through
the equivalence relation of physical indistinguishability. The exclusion
problem is then: if $v$ is supervenient and additive over composition, then
$v$ adds no causal information beyond what the underlying additive physical
valuation already provides. The cocycle picture identifies the locus of
\emph{causal autonomy} of $v$ with the cocycle $\omega_v$: where
$\omega_v(a,b) \ne 0$, the value $v(a\compose b)$ is not determined by
$v(a) + v(b)$ alone, so the joint state $a\compose b$ has a $v$-property
not present in either part.

\paragraph{Status.} The exclusion problem is, when so reformulated,
\emph{circumvented} by any $v$ whose $\omega_v$ is not a coboundary of any
function $f$ supervenient on the physical level. The Lean
\texttt{stronglyEmergent\_iff\_not\_additive} theorem
(Proposition~\ref{prop:strong-iff}) provides the dichotomy: either $v$ is
strongly emergent, in which case the supervenience-respecting cobounding
data $f$ does not exist, or $v$ is additive, in which case Kim's worry is
correct.

\paragraph{Commentary.} The cocycle perspective dissolves the alleged
incoherence of downward causation by relocating it from a metaphysical to a
cohomological problem. ``Downward causation'' is the negation of the
existence of a physically-supervenient $f$ with $\omega_v = \delta f$. Such
non-existence is a perfectly coherent algebraic condition, with examples
\cite{Bedau1997, Anderson1972}. We do not claim that the cocycle
formulation \emph{settles} Kim's exclusion problem in favor of nonreductive
physicalism; we claim only that the problem becomes a tractable algebraic
question once cast in cocycle terms, with the answer depending on the
allowed function class for the cobounding 1-cochain.

\subsection{Chalmers --- strong vs.\ weak emergence}\label{sec:zoo:chalmers}

\paragraph{Original claim.} David Chalmers \cite{ChalmersStrong} distinguishes
\emph{weak emergence} (a high-level property is unexpected given the
low-level facts but in principle deducible from them) from \emph{strong
emergence} (a high-level property is not even in principle deducible from
the low-level facts). Chalmers takes consciousness to be the only credible
candidate for strong emergence in nature; he takes weak emergence to be the
default mode of nontrivial high-level science (\cite{Bedau1997, Holland1998}).
Chalmers's distinction is finer than Anderson's or Mill's because it is
explicitly modal: it concerns what is in-principle derivable, not what is
practically computable.

\paragraph{Formal restatement.} Weak emergence is, again, the bounded
$\omega_v$ regime of Definition~\ref{def:weak}; strong emergence is the
nontriviality of $[\omega_v]$ in $H^2(\IdMon, \mathbb{R})$ relative to the
class of low-level cobounding 1-cochains. The Lean predicate
\texttt{IsStronglyEmergent} is a candidate formalization of the absolute
case; the relative case requires fixing an ``in-principle deducibility''
class and asking whether $\omega_v$ is a coboundary of any $f$ in that
class.

\paragraph{Status.} True (under the chosen translation) by
Theorem~\ref{thm:cocycle} and Proposition~\ref{prop:strong-iff}. The
formalization makes Chalmers's distinction precise:
\begin{itemize}
\item \emph{Bounded $\omega_v$, $\omega_v = \delta f$ for $f$ in the class}:
no emergence, additive valuation.
\item \emph{Bounded $\omega_v$, $\omega_v = \delta f$ for $f$ outside the
class}: weak emergence (Bedau).
\item \emph{Unbounded or non-cobounded $\omega_v$ in the class}: strong
emergence (Chalmers).
\end{itemize}
The Lean lemmas \texttt{boolIndicator\_boundedEmergent} (weak) and
\texttt{boolIndicator\_stronglyEmergent} (strong) demonstrate that both
regimes are realized in a finite, fully verified example.

\paragraph{Commentary.} Chalmers's argument that consciousness is strongly
emergent translates, in the cocycle picture, to the empirical claim that no
physically-cobounding 1-cochain $f$ recovers the cocycle $\omega_v$ of the
``what-it-is-like'' valuation. The cocycle picture is silent on whether
this empirical claim is true; it provides only the vocabulary in which the
question becomes a precise algebraic one. We regard this as substantial
clarificatory progress over the modal phrasing.

\subsection{Holland --- emergence in complex adaptive systems}
\label{sec:zoo:holland}

\paragraph{Original claim.} John Holland's \emph{Emergence: From Chaos to
Order} \cite{Holland1998} develops a theory of emergence in complex adaptive
systems based on agents, interactions, and the cumulative effect of
many-step rule applications. Holland's central technical concept is the
\emph{constrained generating procedure} (CGP): a system of agents whose
interactions, although locally simple, produce globally novel behavior over
many iterations. Examples in Holland's text include cellular automata,
Echo (his agent-based model of artificial economies), classifier
systems, and the games of Go and chess. Emergence, for Holland, is the
build-up of structure that no analysis of single-step rules can predict.

\paragraph{Formal restatement.} Holland's CGP is precisely the chain
$a_1 \compose \cdots \compose a_n$ of Theorem~\ref{thm:telescope}, with
$a_k$ the rule applied at step $k$ and $v$ the valuation recording the
relevant macroscopic observable. The ``cumulative novelty'' of $n$ steps is
$\sum_{k=1}^{n-1} \omega_v(P_k, a_{k+1})$. Holland's claim that emergence
``builds up'' is the assertion that this sum grows nontrivially with $n$ ---
in particular that $\omega_v$ is not a coboundary, since otherwise the
sum would telescope to $f(P_n) - \sum_k f(a_k)$, a difference of
1-cochain values bounded by $\|f\|_\infty$ and hence not growing.

\paragraph{Status.} True, by Theorem~\ref{thm:telescope}, for any chain
$a_1, \dots, a_n$ such that $\omega_v$ accumulates rather than telescopes.
The case in which $\omega_v$ telescopes (i.e.\ is a coboundary) is precisely
the regime of \emph{additive} emergence, where Holland's distinctive
phenomenology disappears. Holland's empirical claim about CGPs is then the
claim that natural CGPs lie in the non-coboundary regime.

\paragraph{Commentary.} Holland's identification of emergence with
\emph{cumulative novelty along chains} is exactly the right intuition for
the cocycle picture, and Theorem~\ref{thm:telescope} is its formal
expression. The Lean lemmas \texttt{comp\_chain\_pair, comp\_chain\_triple,
comp\_chain\_quad, comp\_chain\_pent} of \texttt{Theorems5.lean} unfold the
chain valuation explicitly for $n = 2, 3, 4, 5$, exhibiting the cumulative
$\omega$-sum at each scale. Holland's CGPs sit naturally as the
sub-monoids of $\IdMon$ generated by a finite rule alphabet; the
generating procedure is the chain enumeration, and the constraint is the
choice of admissible $a_k$ at each step.

\subsection{Sawyer --- non-reducible group cognition}\label{sec:zoo:sawyer}

\paragraph{Original claim.} R.\ Keith Sawyer \cite{Sawyer2005} argues that
group-level cognitive phenomena --- collective decision-making, jury
deliberation, jazz improvisation, classroom discourse --- are
\emph{nonreducible} to individual cognition in the sense that no aggregation
function over individuals' mental states predicts the group state. Sawyer
draws on Durkheim's distinction between social facts and individual facts,
on Hutchins's distributed-cognition program, and on his own ethnographic
work on improvisational performance. His thesis is that group cognition has
its own dynamics, irreducible to individual cognition by any aggregation rule.

\paragraph{Formal restatement.} Take $\Ideas$ to be the algebra of
deliberative contributions of group members, $\compose$ the conversational
combination operator, and $v$ a valuation recording group-level cognitive
output (decision, performance, classroom achievement). Sawyer's
nonreducibility is the assertion that no aggregation 1-cochain $f$ on
individual contributions cobounds the group cocycle $\omega_v$, i.e.\ that
$[\omega_v]$ is nontrivial in $H^2$ relative to the class of
``individualistic'' aggregation functions.

\paragraph{Status.} Contingent: Sawyer's nonreducibility is a special case
of Chalmers strong emergence (Section~\ref{sec:zoo:chalmers}), restricted
to the class of aggregation cochains over individuals. The cocycle picture
neither establishes nor refutes Sawyer's claim, but it identifies its
algebraic shape and makes empirical tests well-defined: one tests Sawyer's
thesis by trying, and failing, to construct an $f$ with $\delta f = \omega_v$
within the class of individual-aggregation functions. Failure is a positive
result; success is a refutation.

\paragraph{Commentary.} Sawyer's case is methodologically the cleanest of
the eight: the constraint class on the cobounding cochain is naturally
defined by the social-scientific question (``can we predict group output
from individual inputs alone?''), and the cocycle obstruction is naturally
measured in observed deviations from prediction. The cocycle picture
clarifies that ``group cognition has its own dynamics'' need not contradict
methodological individualism; it requires only that the group-level cocycle
$\omega_v$ have nontrivial class relative to individual-aggregation
1-cochains, while being trivial in the larger class admitting also
interactional coboundary terms. This is the precise sense in which Sawyer's
position is consistent with Holland's: both are non-coboundary statements
about cumulative novelty along chains, differing only in the chosen
constraint class.

\subsection*{Cross-cutting remarks}\label{sec:zoo:cross}

Several of the eight zoo entries make claims that, in the cocycle
vocabulary, become directly comparable. Mill's heteropathy
(Section~\ref{sec:zoo:mill}) is just the predicate $\omega_v \not\equiv 0$,
and so is the negation of additivity. Broad's emergent qualities
(Section~\ref{sec:zoo:broad}) refine this by restricting the class of
admissible cobounding 1-cochains. Anderson's ``more is different''
(Section~\ref{sec:zoo:anderson}) is the additional quantitative claim that
the cocycle term in Theorem~\ref{thm:telescope} dominates the linear sum
asymptotically. Bedau's weak emergence (Section~\ref{sec:zoo:bedau})
sharpens to bounded $\omega_v$, which by Lemma~\ref{lem:sign-tele} prevents
both unbounded super- and unbounded sub-additivity. Kim's exclusion
problem (Section~\ref{sec:zoo:kim}) is the question whether $\omega_v$ is a
coboundary of a supervenience-respecting $f$. Chalmers's strong/weak
distinction (Section~\ref{sec:zoo:chalmers}) is the quotient of these two
predicates --- bounded vs.\ non-cobounded --- and forms a $2\times 2$ grid.
Holland's CGPs (Section~\ref{sec:zoo:holland}) realize the chain
telescoping on a non-coboundary cocycle. Sawyer's group cognition
(Section~\ref{sec:zoo:sawyer}) is Chalmers strong emergence relative to a
specific class of individualistic cobounding functions.

This is not a coincidence. The cocycle picture is what one naturally
arrives at by asking, for any quantitative theory of composite phenomena,
``how does the value of the whole compare to the values of the parts?''
The associativity of composition forces that comparison to live in
$H^2$, and the entire emergentist tradition is the (until-now informal)
exploration of what it means for a $\omega_v$ to be nonzero, bounded,
unbounded, or cohomologically nontrivial in some restricted class.

\section{Closing}

The chapter has assembled the algebraic skeleton of emergence. The single
abstract object is the 2-cocycle $\omega_v$ of a valuation $v$; its
defining identity is associativity-forced (Theorem~\ref{thm:cocycle}); it
is normalized at the unit (Corollary~\ref{cor:normalized-omega}); and it
telescopes along chains (Theorem~\ref{thm:telescope}). The vocabulary of
weak vs.\ strong emergence is recovered from the boundedness vs.\
non-coboundedness of $\omega_v$ in an appropriate cohomology
(Section~\ref{sec:bounded}). Eight historical positions on emergence are
reconstructed as instances of, or constraints on, the cocycle, with
verifiable consequences in each case (Section~\ref{sec:zoo}).

A useful slogan to take from the chapter is the following: in every
quantitative theory of composite phenomena there is a natural valuation $v$,
and \emph{the only thing the cocycle $\omega_v$ records is what we did not
know about the whole from knowing the parts}. The cocycle is therefore not
an exotic addendum to a working theory but the precise residual that
distinguishes the theory from its naively-additive shadow. The size, sign,
boundedness, and cohomology class of $\omega_v$ are then four orthogonal
empirical questions, each independently testable, and each with a
well-defined position in the philosophical literature surveyed in
Section~\ref{sec:zoo}. We have not promised that the cocycle picture
\emph{settles} the disputes between Mill, Broad, Anderson, Bedau, Kim,
Chalmers, Holland, and Sawyer; we have shown only that it furnishes the
common vocabulary in which their disputes become tractable. The
chapters to come exploit that vocabulary at length.

The Lean formalization in \texttt{IdeaTheory.Extensions.EmergenceLevels}
verifies all the technical claims of Sections~\ref{sec:diagnostic}--%
\ref{sec:bounded}: definition of $\omega$, the cocycle identity, both
normalization theorems, the shift formula, the bounded-$\omega$ closure
properties, and the equivalence of strong emergence with non-additivity.
The chain-level telescoping is supported by \texttt{Theorems5.lean}'s
\texttt{CompChain} infrastructure and by \texttt{Theorems7.lean}'s
identity-insertion lemmas, which together verify that the right-hand
side of \eqref{eq:telescope} is well-defined regardless of how the chain is
parenthesized or padded with units.

The next chapter (\cref{ch:grading}) introduces the grading
$\deg : \Ideas \to G$ and uses it to refine the cocycle picture: $\omega_v$
acquires a degree, the telescoping formula respects degrees, and the
weak/strong dichotomy refines into a spectrum indexed by degree. The
chapter after that (\cref{ch:transmission}, on the Transmission Chain)
develops the iterated-composition calculus suggested by
Section~\ref{sec:telescope}, in particular the identity-insertion stability
results of \texttt{Theorems7.lean}.
